\documentclass[11pt,a4paper]{article}
\usepackage[margin=1in]{geometry}
\usepackage{hyperref}
\usepackage{enumitem}
\usepackage{graphicx}
\usepackage{xcolor}

% -----------------------------------------------------
% bvraghav-tiet-ucs503p-article.sty
% -----------------------------------------------------

% Variable: Lab Instructor
\makeatletter \newcommand{\BvrSetLabInstructor}[1]{%
  \gdef\Bvr@LabInstructor{#1}}
% default
\newcommand{\Bvr@LabInstructor}{Dr. Raghav %
  B. Venkataramaiyer}
% public accessor
\newcommand{\BvrLabInstructorValue}{\Bvr@LabInstructor}
\makeatother

% Add subtitle
\usepackage{titling}
% -- Set title bold
\pretitle{\begin{center}\bfseries\fontsize{18bp}{18bp}\selectfont}
\makeatletter
\newcommand{\BvrSubtitle}[1]{%
  \posttitle{%
    \par\end{center}
    \begin{center}\large#1\end{center}
    \vskip 0.5em}%
}
\makeatother

% Hints in the section title(s)
\newcommand{\BvrHint}[1]{%
  {\normalfont\normalsize\itshape\hfill #1}}

% -----------------------------------------------------

\title{Cloth Tracer: Smart Digital Laundry Management System}

\BvrSetLabInstructor{Dr. Raghav B. Venkataramaiyer}

\BvrSubtitle{%
  Project Proposal for UCS503P  \\
  Submitted to: \BvrLabInstructorValue \\ \bfseries
  Thapar Institute of Engineering and Technology% 
}

\author{
  Minkal Goyal, CSED%
  \thanks{Roll No: \texttt{1024030675}, %
    Email: \texttt{mminkal\_be24@thapar.edu}}
  \and 
  Taashvi Aggarwal, CSED%
  \thanks{Roll No: \texttt{1024030681}, %
    Email: \texttt{taggarwal\_be24@thapar.edu}}
\and 
  Anshika Bansal, CSED%
  \thanks{Roll No: \texttt{1024030689}, %
    Email: \texttt{abansal5\_be24@thapar.edu}}
}

\date{\today}

\begin{document}
\maketitle

\section{Higher-order goal%
  \BvrHint{\ldots secondary framing}}
This project aims to improve accountability and complaint resolution in hostel laundry management by supporting the existing laundry process with a digital record system.

The immediate engineering objective is to reduce the difficulty of identifying misplaced clothes, maintaining records and resolving complaints through a simple and practical web-based solution

\section{Time-to-value %
\BvrHint{\ldots primary heuristic}}
\begin{itemize}[leftmargin=0.7cm]
    \item Develop the core student and laundry workflow first and validate it with a small group of users.
    \item Prioritize practical features that can work alongside the existing manual laundry process.
    \item Improve the system based on feedback from students and laundry staff.
\end{itemize}

\section{Problem Statement}
The current hostel laundry process involves large batches of clothes being sorted, washed, dried, ironed and finally sorted according to laundry numbers. The major problems identified are:
\begin{itemize}[leftmargin=0.7cm]
    \item \textbf{Misplacement:} Clothes can sometimes get mixed with another student's laundry.
    \item \textbf{Delayed complaint resolution:} Complaints can take a long time to resolve, especially when the cloth has moved to another hostel.
    \item \textbf{Unidentified clothes:} Laundry staff sometimes have clothes whose laundry numbers are missing or not properly identified.
    \item \textbf{Communication gaps:} Changes in collection dates may not always reach students or caretakers on time.
\end{itemize}

\textbf{Why this matters now (contextual relevance):} A misplaced cloth can require repeated manual searching and coordination between students, laundry staff and caretakers. A digital supporting system can make identification and complaint handling more organized and accountable.

\section{Proposed Solution %
\BvrHint{\ldots Software Proposition}}
\subsection{Overview}
We propose \textbf{ClothTracer} ``Smart Digital Laundry Management System'' system with the following components:
\begin{itemize}[leftmargin=0.7cm]
    \item \textbf{Student Interface:} Upload clothing photos and enter the laundry slip number.
    \item \textbf{Digital Laundry Record:} A record is generated using the submitted slip number and clothing information.
    \item \textbf{Photo-Based Complaints:} Students can raise complaints and attach photos as supporting proof.
    \item \textbf{AI-Assisted Image Matching:} Unidentified clothes can be compared with students' previously uploaded clothing photos to suggest possible matches.
    \item \textbf{Admin/Laundry Interface:} Staff can access relevant records and manage complaints.
     \item \textbf{Collection Notifications:} Changes in laundry collection dates can be communicated through the system.
\end{itemize}

\subsection{Core workflow}

\begin{enumerate}[leftmargin=0.7cm]
    \item Student uploads clothing photos + enters the laundry slip number
    $\rightarrow$ Digital laundry record is generated.

    \item Existing manual slip
    $\rightarrow$ continues to be used for sorting.

    \item Cloth is misplaced/unidentified
    $\rightarrow$ Student raises a complaint + uploads photo as proof.

    \item Unidentified cloth photo
    $\rightarrow$ AI-based image matching
    $\rightarrow$ Possible matches are suggested.

    \item Possible match
    $\rightarrow$ Verified by student/laundry staff
    $\rightarrow$ Complaint is resolved accordingly.
\end{enumerate}

\subsection{Operational constraints}
\begin{itemize}[leftmargin=0.7cm]
    \item The system should support the laundry's existing hostel-wise collection schedule and batch processing.
    
    \item Digital records should be created only with the relevant laundry slip number and clothing details to maintain record reliability.
    
    \item The solution should minimize additional work for laundry staff by allowing students to take and upload photos of their own clothes.
    
    \item Image matching should work as an assistance mechanism, with final identification confirmed by the concerned student or laundry staff.
    
    \item The system should remain usable as the number of students, laundry records and stored images increases.
\end{itemize}

\section{Solution Approach %
\BvrHint{\ldots Engineering focus}}

This is an engineering course project focused on developing a practical digital workflow that supports the existing hostel laundry process and improves record management, accountability and complaint resolution.

\subsection{Digital record and complaint management %
\BvrHint{\ldots workflow support}}
The system will provide:
\begin{itemize}[leftmargin=0.7cm]
    \item Student and staff authentication through Firebase Auth.
    \item Digital laundry records linked with student details and laundry slip numbers.
    \item Complaint submission with relevant evidence for missing or damaged clothes.
    \item Centralized records to support easier complaint investigation and resolution.
\end{itemize}

\subsection{AI-assisted image matching %
\BvrHint{\ldots identification assistance}}
Image matching will be embedding-based:
\begin{itemize}[leftmargin=0.7cm]
    \item Generate image embeddings for stored clothing images using an image embedding model.
    \item Compare the embedding of an unidentified cloth with stored embeddings using image similarity.
    \item Present possible matches for verification by the student or laundry staff.
\end{itemize}

\subsection{Web architecture %
\BvrHint{\ldots web-based solution}}
A typical web-based architecture:
\begin{itemize}[leftmargin=0.7cm]
    \item \textbf{Frontend:} React-based web interface for students and laundry/admin staff.
    \item \textbf{Backend API:} Node.js + Express for authentication integration, laundry records, complaints and application logic.
    \item \textbf{Database:} PostgreSQL for students, staff, slip numbers, laundry records, clothes and complaints.
    \item \textbf{Image Storage:} Cloudinary for clothing, complaint and other relevant images.
\end{itemize}

\subsection{System integration}
The frontend, backend, database, image storage and AI matching components will communicate through APIs, providing a modular architecture that can be developed and deployed independently.

\section{Evaluation Criterion %
\BvrHint{\ldots measurable and attributable}}

The effectiveness of ClothTracer will be evaluated through measurable outcomes related to complaint handling, record accuracy, image matching and usability.

\subsection{Primary evaluation metric}

\textbf{Complaint Resolution Time (CRT):} time taken from complaint registration to identification or resolution of the reported item.

\begin{itemize}[leftmargin=0.7cm]
    \item \textbf{Baseline:} the current manual process may require repeated record checking and, as reported during our laundry interaction, some unresolved cases can extend for weeks or even months.
    \item \textbf{Pilot target:} achieve at least a \textbf{50\% reduction} in the time required to retrieve relevant records and investigate a complaint.
    \item \textbf{Measurement:} compare timestamps and task completion time for selected complaints handled manually and through ClothTracer.
\end{itemize}

\subsection{Secondary evaluation metrics}

\begin{itemize}[leftmargin=0.7cm]

    \item \textbf{Digital Record Accuracy:}
    Target $\geq$ \textbf{90\%} correct association of pilot laundry records with the corresponding student and slip number.

    \item \textbf{AI Matching Effectiveness:}
    Evaluate whether the correct clothing record appears within the system's \textbf{Top-5 suggested matches}; target $\geq$ \textbf{70\%} during pilot testing.

    \item \textbf{Record Retrieval Time:}
    Target retrieval of a relevant laundry record within approximately \textbf{1 minute} under normal pilot conditions.

    \item \textbf{Task Completion Rate:}
    Target $\geq$ \textbf{85\%} successful completion of core tasks such as record creation and complaint submission without assistance.

    \item \textbf{User Usability:}
    Collect student and laundry-staff feedback on a \textbf{5-point scale}, targeting an average score of at least \textbf{4/5}.

\end{itemize}

\subsection{Pilot validation plan %
\BvrHint{\ldots high-level}}

\begin{itemize}[leftmargin=0.7cm]
    \item Conduct an initial pilot with approximately \textbf{10--20 students} and available laundry staff.
    
    \item Create sample laundry records and test representative cases involving normal submissions, complaints and unidentified clothes.
    
    \item Compare manual and ClothTracer-assisted record retrieval and complaint investigation using the defined metrics.
    
    \item Use pilot results and staff feedback to refine the workflow before wider deployment.
\end{itemize}
\section{Scalability %
\BvrHint{\ldots with foresight}}

The system should remain scalable as the number of students, laundry records, complaints and stored images increases.

\begin{enumerate}[leftmargin=0.7cm]

    \item \textbf{Scalable (theoretically):} The system should support increasing data and user load through:
    \begin{itemize}[leftmargin=0.7cm]
        \item separating frontend, backend and AI components,
        \item using indexed database queries for efficient record retrieval,
        \item using cloud storage for scalable image management.
    \end{itemize}

    \item \textbf{Operationally deployable (practically):} The system should support:
    \begin{itemize}[leftmargin=0.7cm]
        \item managed PostgreSQL database,
        \item cloud-based image storage through Cloudinary,
        \item Excel export of laundry and complaint records for administrative use,
        \item platform-based hosting for the web application and backend.
    \end{itemize}

\end{enumerate}

\section{Technology and Engine Availability}

A mature set of technologies and services is available to implement ClothTracer as a web-based system:
\begin{itemize}[leftmargin=0.7cm]
    \item React for the responsive student and admin interface.
    \item Node.js + Express for backend APIs and application logic.
    \item PostgreSQL for structured laundry, student and complaint records.
    \item Firebase Authentication for user authentication and identity management.
    \item Cloudinary for storing clothing and complaint-related images.
    \item Image embedding models for AI-assisted clothing image matching.
\end{itemize}

These established technologies allow the project to focus on practical workflow design, reliable record management and complaint resolution rather than developing infrastructure from scratch.


\section{Project Scope and Deliverables %
\BvrHint{\ldots iteration-friendly}}

\subsection{Initial Deliverables %
\BvrHint{\ldots first iteration}}

\begin{itemize}[leftmargin=0.7cm]
    \item Student and staff authentication.
    \item Student interface for creating digital laundry records.
    \item Laundry slip number-based record generation.
    \item Clothing image storage and retrieval.
    \item Basic complaint submission and management.
    \item Admin/laundry interface for accessing relevant records.
\end{itemize}

\subsection{Subsequent Deliverables}

\begin{itemize}[leftmargin=0.7cm]
    \item AI-assisted matching for unidentified clothing.
    \item Image-based complaint evidence and verification.
    \item Collection-date notifications for students.
    \item Search and filtering of laundry records.
    \item Excel export of laundry and complaint records.
    \item Improved administrative reports and record management.
\end{itemize}


\section{Risks and Mitigations}

\begin{itemize}[leftmargin=0.7cm]
    \item \textbf{Data privacy:} store only required student and clothing information and restrict access to authorized users.
    
    \item \textbf{Staff adoption:} keep the workflow simple and retain the existing manual slips for sorting.
    
    \item \textbf{Incorrect AI matches:} present AI results as possible matches requiring human verification rather than automatic identification.
    
    \item \textbf{Image quality:} provide basic guidelines for capturing clear clothing images.
    
    \item \textbf{Operational constraints:} design the system to support the existing batch-based laundry process without requiring major changes to the current workflow.
\end{itemize}


\section{Summary}

ClothTracer proposes a practical web-based laundry management system to improve record reliability, accountability and complaint resolution in hostel laundry operations. The system supports the existing batch-based process through digital records, clothing images, structured complaint management and AI-assisted identification of unidentified clothes. It focuses on practical engineering, usability and scalability while keeping the existing laundry workflow and manual sorting process intact.
\end{document}
